\documentclass{rvcepaper}
% s7-common.tex -- macros shared by the Sem 7 CIE papers (\input right after \documentclass{rvcepaper})
% placeholder for text that cannot be read in the scan
\newcommand{\illeg}{\textit{[illegible]}}
\newcommand{\illegc}{{\tiny\textit{[illegible]}}}
% question table with a custom header row: \begin{gtable}[PART -- A]{Sl.\newline No.}{Questions}{M}{BT}{CO}
% the optional title is set as the table's first row so it never strands at a page foot
\newcommand{\gtitlerow}{}
\newenvironment{gtable}[6][]{%
  \def\gtitle{#1}\ifx\gtitle\empty\def\gtitlerow{}\else\def\gtitlerow{\multicolumn{5}{c}{\bfseries\gtitle}\\*[2pt]}\fi
  \begin{longtable}{|C{1.05cm}|p{12.2cm}|C{0.95cm}|C{0.95cm}|C{0.95cm}|}
  \gtitlerow
  \hline \textbf{#2} & \multicolumn{1}{c|}{\textbf{#3}} & \textbf{#4} & \textbf{#5} & \textbf{#6}\\* \hline
  \endfirsthead
  \endhead
}{\end{longtable}}
\newcommand{\twonotes}[2]{{\small Note:}\par\vspace{-1pt}
\hspace*{0.5cm}{\footnotesize 1. \hspace{0.1cm}#1}\par
\hspace*{0.5cm}{\footnotesize 2. \hspace{0.1cm}#2}\par\vspace{3pt}}
\newcommand{\onenote}[1]{{\small Note:}\par\vspace{-1pt}
\hspace*{0.5cm}{\footnotesize 1. \hspace{0.1cm}#1}\par\vspace{3pt}}
% keep the BT/CO caption and the marks-distribution table on the same page
\renewcommand{\btfoot}[2]{\par\vspace{8pt}\noindent\begin{minipage}{\textwidth}\centering
  {\small BT-Blooms Taxonomy, CO-Course Outcomes, M-Marks}\par\vspace{2pt}
  \begin{tabular}{#1}#2\end{tabular}\end{minipage}\par\nopagebreak\vspace{2pt}}
\newcommand{\xfoot}[3]{\par\vspace{8pt}\noindent\begin{minipage}{\textwidth}\centering
  {\small #1}\par\vspace{2pt}\begin{tabular}{#2}#3\end{tabular}\end{minipage}\par\nopagebreak\vspace{2pt}}
\newcommand{\qq}[5]{#1 & #2 & #3 & #4 & #5\\}
\newcommand{\qm}[3]{#1 & #2 & #3\\ \hline}
% other-department header table (Date / Course Code / Sem rows + title row)
\newcommand{\ofour}[7]{\noindent\begingroup\renewcommand{\arraystretch}{1.3}\fontsize{11.5}{13}\selectfont
\begin{tabular}{|C{3.6cm}|C{4.6cm}|C{4.0cm}|C{3.8cm}|}\hline
Date & #1 & Maximum Marks & #2\\\hline
Course Code & #3 & Duration & #4\\\hline
Sem & #5 & \multicolumn{2}{c|}{#6}\\\hline
\multicolumn{4}{|c|}{\bfseries\fontsize{13}{15}\selectfont #7}\\\hline
\end{tabular}\endgroup\par\vspace{6pt}}
\renewcommand{\arraystretch}{1.05}
% the closing stars line stays with the block above it
\renewcommand{\stars}[1]{\par\nopagebreak\vspace{2pt}\nopagebreak\centerline{#1}}
% printed value struck through by hand, with the handwritten correction under it: \corr{60}{34}
\usepackage[normalem]{ulem}
\newcommand{\corr}[2]{\shortstack{\sout{#1}\\#2}}
\usetikzlibrary{decorations.pathmorphing}
% "Academic year" header table without the DEPARTMENT OF block (use inside iempaperB after \paperinfo)
\newcommand{\iemtableB}{%
  \par\vspace{4pt}\noindent\begingroup\renewcommand{\arraystretch}{1.5}\bfseries\fontsize{11.5}{13}\selectfont
  \begin{tabular}{|p{4.2cm}|p{5.0cm}|p{3.7cm}|p{3.6cm}|}\hline
   Date & \getkey{date} & Maximum Marks & \getkey{maxmarks}\\\hline
   Course Code & \getkey{code} & Duration & \getkey{duration}\\\hline
   Sem & \getkey{sem} & \multicolumn{2}{l|}{\getkey{exam}}\\\hline
   \multicolumn{4}{|c|}{\getkey{title}}\\\hline
  \end{tabular}\endgroup\par\vspace{4pt}}
% other-department CIE header: \xdept{Department line}{left1}{mid1}{right1}{left2}{mid2}{right2}{course title cell}{right3}
\newcommand{\xdept}[9]{\par\vspace{2pt}\centerline{\bfseries\fontsize{13.5}{16}\selectfont #1}\par\vspace{4pt}
  \noindent\begingroup\renewcommand{\arraystretch}{1.5}\bfseries\fontsize{11.5}{13}\selectfont
  \begin{tabular}{|p{4.8cm}|C{5.6cm}|p{5.7cm}|}\hline
  #2 & #3 & #4\\\hline #5 & #6 & #7\\\hline \multicolumn{2}{|p{10.8cm}|}{#8} & #9\\\hline
  \end{tabular}\endgroup\par\vspace{8pt}}

\begin{document}

% ---------- HS271TA Indian Knowledge System (Improvement CIE)
\begin{iempaperB}
\paperinfo{ayear={2025-2026 (Odd Semester)},usn=0,date={Jan 2026},code={HS271TA},sem={VII},exam={\hfill Improvement CIE\hfill\mbox{}},maxmarks={50 + 10},duration={120 Min},
 title={INDIAN KNOWLEDGE SYSTEM},shownote=0}
\iemtableB
\begin{gtable}[{\hfill PART -- A\hfill\llap{\normalfont(10 Marks)}}]{Sl.No}{}{M}{CO}{BT}
\q{1}{Which value links IKS with SDGs?\newline A. Consumerism \qquad B. Sustainability\newline C. Militarism \qquad D. Competition}{1}{3}{1}
\q{2}{Janapada in Arthaśāstra refers to:\newline A. Army \quad B. Territory and people \quad C. Treasury \quad D. Allies}{1}{3}{1}
\q{3}{SDGs are adopted by:\newline A. IMF \quad B. World Bank \quad C. United Nations \quad D. WTO}{1}{3}{2}
\q{4}{The king's happiness lies in:\newline A. Luxury \quad B. Conquest \quad C. Welfare of subjects \quad D. Fame}{1}{2}{1}
\q{5}{Sustainable governance means:\newline A. Short-term profits \qquad B. Resource exploitation\newline C. Intergenerational equity \qquad D. Militarization}{1}{2}{1}
\q{6}{Governance in the Rāmāyaṇa is best described as:\newline A. Autocratic \qquad B. Tyrannical \qquad C. Welfare-oriented \qquad D. Laissez-faire}{1}{3}{1}
\q{7}{Arthaśāstra views the king as:\newline A. Divine \qquad B. Absolute ruler \qquad C. Servant of the state \qquad D. Merchant}{1}{2}{1}
\q{8}{The idea of ``Lokasaṅgraha'' means:\newline A. Individual gain \qquad B. Social welfare\newline C. Military power \qquad D. Taxation}{1}{3}{1}
\q{9}{Which ancient text integrates diplomacy, war and peace?\newline A. Upaniṣads \quad B. Rāmāyaṇa \quad C. Arthaśāstra \quad D. Vedas}{1}{3}{1}
\q{10}{Which is NOT one of the Saptaṅga theory elements?\newline A. Swāmin \quad B. Mantri \quad C. Dharma \quad D. Janapada}{1}{2}{1}
\end{gtable}
\begin{longtable}{|C{1.05cm}|p{12.2cm}|C{0.95cm}|C{0.95cm}|C{0.95cm}|}
\multicolumn{5}{c}{\bfseries\hfill TEST\hfill\llap{\normalfont(50 marks)}}\\*[2pt]
\hline \textbf{Sl.\newline No.} & \multicolumn{1}{c|}{\textbf{Part- B Questions}} & \textbf{M} & \textbf{CO} & \textbf{BT}\\* \hline
1 a & Define Rāja Dharma. & 02 & \multirow{4}{*}{4} & 1\\*
b & What are the key governance attributes highlighted in the Rāmāyaṇa? & 02 & & 2\\*
c & Conclude the relevance of ancient Indian governance. & 03 & & 2\\*
d & Explain Saptaṅga theory of state. & 03 & & 2\\\hline
2 a & Describe governance principles from Rāma's rule. & 03 & \multirow{4}{*}{3} & 2\\*
b & Compare ancient and modern governance values. & 03 & & 2\\*
c & Why is Kosa considered vital for administration? & 02 & & 1\\*
d & Name the wives of Bharath and Shathrghan from Ramayana & 02 & & 2\\\hline
3 a & Explain the ideal ruler concept in the Rāmāyaṇa. & 03 & \multirow{4}{*}{3} & 1\\*
b & How does Arthaśāstra emphasize accountability? & 03 & & 2\\*
c & What role did ethics play in administration? & 02 & & 1\\*
d & What are SDGs? & 02 & & 1\\\hline
4 a & Mention the Principles of Effective Governance in the Arthaśāstra & 03 & \multirow{4}{*}{4} & 2\\*
b & Explain the Interrelationship of the Concepts of arthashatra & 03 & & 2\\*
c & Mention any two Governance Principles Illustrated in Rāmāyaṇa & 02 & & 1\\*
d & Name any two SDGs with the theme & 02 & & 1\\\hline
5 a & What is sustainable governance? & 02 & \multirow{4}{*}{4} & 1\\*
b & What lessons does Rāmāyaṇa offer modern administrators? & 03 & & 1\\*
c & How does Dharma relate to SDGs? & 03 & & 2\\*
d & What does Durga signify in governance? & 02 & & 2\\\hline
\end{longtable}
\par\vspace{2pt}\centerline{\small\itshape [The scan of this page stops at Q5 d; anything after it is not in the scanned copy.]}
\end{iempaperB}

% ---------- IM372IA Product Design and Development (Improvement CIE)
\begin{iempaper}
\paperinfo{shownote=0,date={12th January 2026},exam={Improvement CIE},maxmarks={10 + 50},sem={7\textsuperscript{th}},prog={UG},duration={30 + 90 Min},
 title={Product Design and Development},code={IM372IA}}
\iemhead
\onenote{Answer all the Questions.}
\begin{cietable}{M}{BT}{CO}
\qpart{Part -- A}
\q{1.}{Give one example of poor human-machine interaction.}{01}{L$_1$}{CO2}
\q{2.}{How does AI support cognitive ergonomics?}{01}{L$_2$}{CO3}
\q{3.}{What is workspace optimization using AI?}{01}{L$_2$}{CO4}
\q{4.}{What is the main goal of DFM?}{01}{L$_2$}{CO3}
\q{5.}{What is meant by the approach to design for manufacturing?}{01}{L$_2$}{CO4}
\q{6.}{State one principle of prototyping.}{01}{L$_1$}{CO4}
\q{7.}{What is the first step in planning for prototypes?}{01}{L$_2$}{CO4}
\q{8.}{What is virtual prototyping?}{01}{L$_2$}{CO4}
\q{9.}{What is concurrent engineering?}{01}{L$_1$}{CO4}
\q{10.}{Give one advantage of virtual prototyping over physical prototyping.}{01}{L$_2$}{CO3}
\qpart{Part -- B}
\q{1.}{Consider any of the following products and well-known brand related to kitchen supplies, tools, office supply, gifts and a set of competing products. (For Example, a Traditional Cork Screw Bottle opener and a Big Wing Corkscrew Bottle opener). Assess both the brand and competing product on five ID quality categories such as Quality of the user interaction, Emotion Appeal, Ability to Maintain and Repair the Product, Appropriate use of resources and Product Differentiation. Use levels of Low, Medium and High for the assessment to identify the Level of Importance in each of the five categories. Which product would you purchase? Would you be willing to pay more for it than for the others?}{10}{L$_4$}{CO4}
\q{2.}{What are the four purposes for which prototyping can be used? Discuss each of these types briefly. A furniture manufacturer was considering a line of seating products to be fabricated by cutting and bending a recycled plastic product available in large sheets. Create a Sketch of the prototype of at least one possible chair design assuming paper or cardboard as prototyping material. What can you learn about the chair design?}{10}{L$_4$}{CO3}
\q{3}{What are the benefits of developing the physical prototypes during design and development of procts? Identify and position the following prototypes
\begin{romi}\item[i.] Boeing wing prototype\item[ii.] Bahco screwdrivers\item[iii.] Projector\item[iv.] Bicycle in SolidWorks\item[v.] Urethane casting\end{romi}
in the plot along the four quadrants along X axis (Ranging from Focussed to Comprehesive) along Y axis (Rannging from physical to analytical)Briefly expl\illeg\ your reasoning for each placement.}{10}{L$_4$}{CO3}
\q{4.}{Explain the differences between a ``works-like''ing prototype and a ``looks- alike ``prototype with a sketch of both for a product like a Nail cutter or any other consumer product of your choice. (Showing the detailing in the sketch is absolutely mandatory)\newline
Many product development teams separate the ``looks-like'' prototype from the ``works-like'' prototype. They do this because integrating both function and form is difficult in the early phases of development. What are the strengths and weaknesses of this approach? For what types of products might this approach be dangerous?}{05}{L$_4$}{\illegc}
\q{5a.}{What are the principles of human-machine interaction, and how they are applied in a medical diagnostic kiosk?}{05}{L$_4$}{\illegc}
\q{b.}{How is ergonomics integrated into IoT-based smart kitchen appliances?}{05}{L$_4$}{\illegc}
\end{cietable}
\end{iempaper}

% ---------- IM373TA Total Quality Management (Improvement CIE)
\begin{iempaper}
\paperinfo{shownote=0,date={13/01/2026},exam={Improvement CIE},maxmarks={10 + 50},sem={VII},prog={UG},duration={30 + 90 Min},
 title={Total Quality Management},code={IM373TA}}
\iemhead
{\small Note: Answer all questions}\par\vspace{3pt}
\begin{cietable}{M}{BT}{CO}
\qpart{Part -- A}
\q{1.}{Define DPMO.}{2}{1}{CO1}
\q{2.}{What is the primary goal of Six Sigma?}{2}{1}{CO3}
\q{3.}{What is the role of a Champion in Six Sigma?}{2}{2}{CO2}
\q{4.}{List the five phases of the standard Six Sigma improvement methodology.}{2}{2}{CO3}
\q{5.}{Why is verification important in DFSS?}{2}{3}{CO4}
\qpart{Part -- B}
\q{1}{Explain the historical development of Six Sigma.}{10}{2}{CO2}
\q{2}{Who are the key players in a Six Sigma project and what are their roles? Enumerate.}{10}{3}{CO2}
\q{3}{Explain how Failure Modes and Effects Analysis (FMEA) supports Design for Six Sigma implementation.}{10}{2}{CO2}
\q{4}{Choose a product and explain how do you apply Design for Six Sigma in product development?}{10}{2}{CO4}
\q{5}{Explain how Design for Six Sigma can be applied to service industries.}{10}{2}{CO4}
\end{cietable}
\end{iempaper}

% ---------- IM374TFB Artificial Intelligence in Industrial and Systems Engineering (Improvement CIE)
\begin{iempaper}
\paperinfo{shownote=0,date={12\textsuperscript{th} JAN 2026},exam={IMPROVEMENT CIE},maxmarks={10 + 50},sem={VII},prog={UG},duration={30 + 90 Min},
 title={Artificial Intelligence in Industrial and Systems Engineering},code={IM374TFB}}
\iemhead
\onenote{Answer all the Questions.}
\begin{cietable}{M}{BT}{CO}
\qpart{Part -- A}
\q{1}{Define AI-based ergonomic risk assessment.}{2}{3}{02}
\q{2}{What is pose estimation in human motion analysis?}{2}{2}{01}
\q{3}{Mention two benefits of human--robot collaboration.}{2}{2}{03}
\q{4}{What is ergonomic risk assessment using AI?}{2}{3}{01}
\q{5}{Define human--robot collaboration (HRC).}{2}{2}{02}
\qpart{Part -- B}
\q{1}{Explain intelligent assistance systems and their role in supporting industrial workers.}{10}{3}{01}
\q{2}{How does AI help in evaluating task performance in industrial environments?}{10}{3}{02}
\q{3}{Explain how AI improves safety in human--robot collaboration.}{10}{3}{03}
\q{4}{Explain how AI-powered human motion analysis improves ergonomic evaluation.}{10}{4}{03}
\q{5}{Discuss the role of computer vision in ergonomic risk assessment.}{10}{3}{03}
\end{cietable}
\end{iempaper}

% ---------- IM374TFA Predictive Analytics (CIE - Improvement)
\begin{iempaper}
\paperinfo{shownote=0,date={12\textsuperscript{th} Jan 2026},exam={CIE -- (Improvement)},maxmarks={10 + 50},sem={7\textsuperscript{th}},prog={UG},duration={30 + 90 Min},
 title={Predictive Analytics},code={IM374TFA}}
\iemhead
\onenote{Answer all the Questions.}
\begin{cietable}{M}{BT}{CO}
\qpart{Part -- A}
\q{1.}{An industrial engineering firm is using sensor data from its machinery to predict time-to-failure (a regression problem). The data contains thousands of highly correlated features and noise. Which ensemble method is generally preferred for achieving high predictive accuracy in complex, noisy, non-linear industrial applications?}{02}{L$_3$}{CO3}
\q{2.}{A financial management firm uses an ensemble model to classify loan applications as high-risk or low-risk. Due to strict regulatory requirements, they need to ensure the final model is highly robust and its prediction variance is minimized, even if individual decision trees are slightly biased. Which ensemble method is inherently designed for variance reduction through parallel training on bootstrapped data?}{02}{L$_3$}{CO3}
\q{3.}{What is the goal of Principal Component Analysis?}{01}{L$_1$}{CO3}
\q{4.}{Anomaly detection aims to identify what in the Data pattern?}{01}{L$_2$}{CO4}
\q{5.}{Which data mining technique is primarily used to discover relationships between items in large transaction datasets, such as finding products frequently bought together in a retail setting?}{01}{L$_2$}{CO3}
\q{6.}{The first common first step in integrating a data mining process into an existing industrial workflow, after defining the problem, is to assess the quality of what and availability of what with respect to relevance to industrial scenario?}{01}{L$_1$}{CO4}
\q{7.}{A company wants to predict which machine in a factory is most likely to fail in the next 7 days based on sensor data. What is the most appropriate \textit{supervised} data mining technique for this \textbf{predictive maintenance} problem?}{02}{L$_4$}{CO4}
\qpart{Part -- B}
\q{1a.}{Covariance matrix: $\Sigma = \begin{matrix}4&2&0\\2&3&0\\0&0&1\end{matrix}$\par\vspace{2pt}
Compute eigenvalues and proportion of variance explained by the first two principal components.}{10}{L$_4$}{CO3}
\q{2a.}{Throughputs: ([95, 100, 98, 97, 96, 150, 99, 101]). Compute (Q1, Q3), IQR, and flag outliers using (x $<$ Q1 - 1.5* *mathrm\{IQR\}) or (x $>$ Q3 + 1.5*cdot *mathrm\{IQR\}).}{06}{L$_3$}{CO4}
\q{b.}{A company is analyzing customer feedback from Twitter. Which specific sub-field of \textbf{Text Mining} would be most directly applied to automatically classify these tweets as positive, negative, or neutral?}{04}{L$_4$}{CO3}
\q{3}{Critically evaluate the statement: ``For any supervised learning problem, if a Random Forest and a Gradient Boosting Machine are both meticulously tuned, the Gradient Boosting Machine will theoretically always achieve superior performance.'' Your evaluation must discuss the inherent model structure and specific data conditions (e.g., noise, outliers) where this statement might fail.}{10}{L$_4$}{CO4}
\q{4a.}{In a financial fraud detection case study, a high volume of false positives (labeling legitimate transactions as fraudulent) is observed. Propose one simple strategy to adjust the model or its outputs to reduce the false positive rate.}{05}{\illegc}{\illegc}
\q{b.}{Assess the primary ethical consideration raised when using credit card transaction data (even if anonymized) in a finance industry case study for building a predictive model of spending habits.}{05}{\illegc}{\illegc}
\q{5a.}{Design and justify the most appropriate data mining technique (classification, clustering, or association rule mining) for a healthcare insurance company aiming to identify groups of fraudulent claim patterns that they have never seen before.}{06}{\illegc}{\illegc}
\q{b.}{You are integrating a demand forecasting model (regression) into a supply chain's existing inventory management system. Justify the need for a rigorous A/B testing phase before full deployment of the new model.}{04}{\illegc}{\illegc}
\end{cietable}
\end{iempaper}

% ---------- IM374TFD Sustainability Engineering & Circular Economy (Improvement CIE; date printed "12 January 2025")
\begin{iempaper}
\paperinfo{shownote=0,date={12\textsuperscript{th} January 2025},exam={Improvement CIE},maxmarks={10 + 50},sem={7\textsuperscript{th}},prog={UG},duration={30 + 90 Min},
 title={Sustainability Engineering \& Circular Economy},code={IM374TFD}}
\iemhead
\onenote{Answer all the Questions.}
\begin{gtable}[PART -- A]{Sl.\newline No.}{Questions}{M}{BT}{CO}
\q{1}{State two economic advantages of Circular Economy.}{2}{2}{4}
\q{2}{What does the 7R model stand for?}{2}{2}{4}
\q{3}{What is the significance of ``Design for Disassembly''?}{2}{2}{4}
\q{4}{Define resource efficiency in the context of Circular Economy.}{2}{1}{4}
\q{5}{What is the link between Circular Economy and SDG 12?}{2}{2}{1}
\end{gtable}
\begin{gtable}[PART -- B]{Sl.\newline No.}{Questions}{M}{BT}{CO}
\q{1.}{Discuss the integration of Sustainable Engineering with Circular Economy.}{10}{2}{1}
\q{2.}{Describe the framework and design strategies of Circular Economy.}{10}{2}{4}
\q{3.}{Propose a Circular Eco-Design for a common product- Smart phone.}{10}{5}{4}
\q{4.}{Discuss the role of Circular Economy in shaping future environmental governance.}{10}{3}{4}
\q{5.}{Propose a Circular Economy strategy for garment manufacturing sector.}{10}{5}{4}
\end{gtable}
\stars{***********}
\end{iempaper}

% ---------- IM375TGN Global Supply Chain Management (CIE - Improvement)
\begin{iempaper}
\paperinfo{shownote=0,date={14/01/2026},exam={CIE -- \ Improvement},maxmarks={10 + 50},sem={VII},prog={UG},duration={30 + 90 Min},
 title={Global Supply Chain Management},code={IM375TGN}}
\iemhead
\onenote{Answer all the Questions.}
\begin{cietable}{M}{BT}{CO}
\qpart{Part -- A}
\q{1.}{A smartphone manufacturer decides to procure chargers from an external vendor instead of producing them internally. This is an example of a \blank[1.4cm] decision.}{1}{2}{3}
\q{2.}{In supplier selection, \blank[1.4cm] are used to create competition among suppliers to achieve lower prices.}{1}{2}{3}
\q{3.}{Revenue management is especially important for \blank[1.4cm] assets such as airline seats and hotel rooms.}{1}{1}{1}
\q{4.}{Vendor-Managed Inventory (VMI) is a replenishment approach where the \blank[1.4cm] is responsible for managing inventory at the customer location.}{1}{2}{3}
\q{5.}{Sharing risk and reward among supply chain partners helps improve \blank[1.4cm] and overall performance.}{1}{1}{3}
\q{6.}{A manufacturer experiences large order fluctuations despite stable customer demand due to poor information sharing. This phenomenon is called the \blank[1.4cm] effect.}{1}{1}{1}
\q{7.}{A retail company hires an external firm to manage transportation, warehousing, and order fulfilment. This firm is best described as a \blank[1.4cm] provider.}{1}{2}{3}
\q{8.}{A supermarket allows its supplier to monitor sales data and replenish stock automatically to avoid stockouts. This system is known as \blank[1.4cm].}{1}{2}{3}
\q{9.}{An airline increases ticket prices as the departure date approaches due to limited seat availability. This practice illustrates \blank[1.4cm] pricing.}{1}{2}{1}
\q{10.}{A buyer and supplier jointly agree to share profits from cost savings achieved through collaboration. This represents \blank[1.4cm] sharing in the supply chain.}{1}{1}{3}
\qpart{Part -- B}
\q{1}{Compare Vendor-Managed Inventory (VMI) and Collaborative Planning, Forecasting, and Replenishment (CPFR), highlighting their roles in achieving supply chain coordination.}{10}{3}{3}
\q{2}{Describe the Supply Chain IT framework and explain how the integration of planning, execution, and analytics systems helps organizations achieve competitive advantage.}{10}{2}{4}
\q{3}{Explain the bullwhip effect in supply chains and critically discuss comprehensive strategies adopted to minimize its impact.}{10}{3}{1}
\q{4}{Elaborate on the make-or-buy decision in sourcing and develop a decision framework incorporating strategic, financial, and operational considerations.}{10}{4}{3}
\q{5}{Explain pricing and revenue management strategies in supply chains and discuss their role in balancing demand and supply.}{10}{2}{3}
\end{cietable}
\end{iempaper}

% ---------- CV375TGG Solid Waste Management & Statuatory Rules (Test-3, Improvement Test)
\begin{iempaperB}
\paperinfo{ayear={2025-2026 (Odd Sem)},usn=0,dept={CIVIL ENGINEERING},date={14.1.2026},code={CV375TGG},sem={VII},exam={Test-3 (Improvement Test)},maxmarks={60},duration={90+20 Min},
 title={SOLID WASTE MANAGEMENT \& STATUATORY RULES},shownote=0}
\iemheadB
\begin{gtable}[PART A]{Q. No}{Question}{Marks}{CO}{BTL}
\q{1}{Define E-waste.}{2}{3}{1}
\q{2.}{Minimum stack height for bio-medical waste incinerator should be \blank[0.8cm] mtrs above ground.}{1}{4}{1}
\q{3}{The thickness of carry bags according to 2016 amendment for plastic waste (management \& handling ) rules, is \blank[1.2cm]}{1}{4}{2}
\q{4}{E-waste (management \& handling) rules, was promulgated in the year \blank[0.8cm]}{1}{3}{1}
\q{5}{Define Extended producer responsibility as per E- waste rules.}{1}{3}{2}
\q{6}{The temperature of primary chamber in the biomedical waste incinerator shall be \blank[1.2cm]}{1}{4}{1}
\q{7}{Schedule \blank[0.8cm] of Bio-medical waste management rules corresponds to label for transport of Bio-medical waste containers or bags.}{1}{4}{1}
\q{10}{Enumerate the equation to calculate the combustion efficiency of the biomedical waste incinerator.}{1}{4}{1}
\q{11}{What are multi layered plastics.}{1}{3}{1}
\end{gtable}
\begin{gtable}[PART B]{Q.\newline No.}{Questions}{Marks}{\scriptsize Course Outcomes}{BTL}
\q{1}{Explain briefly different categories of Bio-medical waste and their treatment/disposal options.}{10}{4}{1}
\q{2}{Enumerate the salient features of E-waste (management \& handling) rules, 2016.}{10}{3}{1}
\qq{3}{a. \ State the adverse effects of plastic waste.}{5}{3}{2}
\q{}{b. \ Illustrate the standards for deep burial of bio- medical waste.}{5}{4}{1}
\q{4}{Explain the roles and responsibilities of E-waste producer, recycler and dismantler.}{10}{4}{2}
\q{5}{Discuss the salient features of plastic waste (management \& handling) rules, 2016 with amendments.}{10}{4}{\illegc}
\end{gtable}
\stars{***********}
\end{iempaperB}

% ---------- HS375TGT Principles & Practices of Cyber Law (Improvement CIE)
\begin{iempaper}
\paperinfo{shownote=0,date={14 Jan 2026},exam={Improvement CIE},maxmarks={10 + 50},sem={VII},prog={UG},duration={30 + 90 Min},
 title={Principles \& Practices Of \ Cyber Law},code={HS375TGT}}
\iemhead
\onenote{Answer all the Questions.}
\begin{cietable}{M}{BT}{CO}
\qpart{Part -- A}
\q{1.}{What is meant by ``Right to be Forgotten'' under GDPR?}{2}{2}{2}
\q{2.}{Define Data Breach?}{2}{2}{2}
\q{3.}{What is the main objective of PIPEDA?}{2}{2}{2}
\q{4.}{List any \ two \ challenges to intellectual property (IP) protection in cyberspace?}{2}{2}{4}
\q{5.}{An E-Commerce site uses a software method to classify products based on customer behavior. Why will the Indian Patent Office treat this as a business method?}{2}{2}{4}
\qpart{Part -- B}
\q{1}{Analyse the rights granted to data subjects under the GDPR. How do these rights strengthen informational privacy and digital autonomy?}{10}{3}{2}
\q{2}{A hospital uses a cloud-based system to store patient records, including biometrics, medical history, Aadhar details, and insurance information. A misconfigured server allowed public access for several days before discovery. Explain the types of data exposed in this incident. Analyze why protecting such data is critical in cyberspace and discuss the legal liabilities under the IT Act and the Digital Personal Data Protection Act, 2023.}{10}{3}{2}
\q{3}{A website hosts thousands of pirated e-books and music files without permission from copyright owners. The domain name closely resembles a famous publisher's trademark, and the site also uses a customized algorithm to hide copyright-protected content. Discuss the IP issues involved and identify legal remedies available under Indian law.}{10}{4}{4}
\q{4}{A mobile app mimics the interface and services of a competing e-commerce platform. It copies slogans, product images, and uses a domain name with minor spelling variation. Users are tricked into believing it is the official app. Analyse the copyright, trademark, and unfair competition issues. Suggest legal remedies.}{10}{4}{4}
\q{5}{Discuss the major trademark issues that arise in cyberspace. Explain the distinction between domain names and trademarks, the nature of domain name disputes and the relevant legal framework governing them.}{10}{3}{4}
\end{cietable}
\end{iempaper}

% ---------- ME375TGP Industry 5.0 (Institutional Elective), Improvement CIE
\begin{xpaper}
\par\vspace{4pt}
\noindent\begingroup\renewcommand{\arraystretch}{1.3}\bfseries\fontsize{11.5}{13}\selectfont
\begin{tabular}{|C{2.4cm}|C{5.4cm}|C{2.3cm}|C{2.3cm}|C{2.3cm}|}\hline
\multicolumn{5}{|c|}{DEPARTMENT OF MECHANICAL ENGINEERING}\\
\multicolumn{5}{|c|}{AY 2025-26 (ODD SEM)}\\\hline
Date & \shortstack{13/01/2026\\Tuesday 1.20 PM to 3.20 PM} & \shortstack{Improvement\\CIE} & Max. Marks & 10+50\\\hline
Semester & VII & UG & Duration & 110 mins\\\hline
Course & \multicolumn{2}{c|}{INDUSTRY 5.0 (Institutional Elective)} & Code & ME375TGP\\\hline
\end{tabular}\endgroup\par\vspace{10pt}
\begin{gtable}{\#}{QUIZ}{M}{BT}{CO}
\q{\textbf{1}}{What role do cobots (collaborative robots) play in Industry 5.0?}{1}{1}{1}
\q{\textbf{2}}{How does Industry 5.0 promote sustainability?}{1}{1}{1}
\q{\textbf{3}}{Which technology enables personalized manufacturing in Industry 5.0?}{1}{2}{2}
\q{\textbf{4}}{A potential challenge associated with adopting Industry 5.0 is ----------}{1}{2}{2}
\q{\textbf{5}}{How does Industry 5.0 redefine the role of humans in manufacturing compared to Industry 4.0?}{1}{2}{2}
\q{\textbf{6}}{In the context of Industry 5.0, what does ``mass personalization'' refer to?}{1}{1}{1}
\q{\textbf{7}}{A critical step in the Industry 5.0 transition roadmap is workforce development. What does this primarily involve?}{1}{2}{2}
\q{\textbf{8}}{Why is energy-efficient and ``green'' AI a key research direction for Industry 5.0?}{1}{1}{1}
\q{\textbf{9}}{What is the purpose of a ``digital twin'' in the context of Industry 5.0 research and trends?}{1}{2}{2}
\q{\textbf{10}}{In the aerospace sector, Boeing has successfully trialed which Industry 5.0-aligned technology to reduce worker fatigue and increase productivity?}{1}{2}{2}
\end{gtable}
\par\vspace{10pt}
\begin{gtable}{\#}{TEST}{M}{BT}{CO}
\q{\textbf{1}}{Bring out the important differences between the Industry 4.0 \& Industry 5.0 frameworks. Also mention the benefits of the transition from Industry 4.0 to 5.0.}{10}{3}{1}
\q{\textbf{2}}{Elaborate the key characteristics, technologies \& benefits of Industry 4.0.}{10}{4}{2}
\q{\textbf{3}}{Name the three core pillars of Industry 5.0 and explain the key technologies actions involved.}{10}{2}{2}
\q{\textbf{4}}{Illustrate some case studies of successful implementation of Industry 5.0 from automotive, aerospace, and electronics industries.}{10}{3}{2}
\q{\textbf{5}}{Summarize the research directions that are ongoing to integrate industrial innovation with broader societal values.}{10}{4}{3}
\end{gtable}
\stars{********}
\end{xpaper}

% ---------- EI375TGK Disease and Diagnostics an Engineering Perspective (Improvement CIE)
\begin{iempaperB}
\paperinfo{ayear={2025-2026(ODD Sem)},usn=0,dept={ELECTRONICS AND INSTRUMENTATION ENGINEERING},date={14.01.2026},code={EI375TGK},sem={VII},exam={Improvement CIE},maxmarks={10+50},duration={120 Min},
 title={Course Title: DISEASE AND DIAGNOSTICS AN ENGINEERING PERSPECTIVE},shownote=0}
\iemheadB
\begin{gtable}{Sl. No.}{PART-A}{M}{BT}{CO}
\q{1.1}{What is Artificial Intelligence (AI) in healthcare?}{2}{L1}{4}
\q{1.2}{Give one example of structured data in healthcare.}{1}{L2}{4}
\q{1.3}{Give one example of unstructured healthcare data.}{1}{L2}{4}
\q{1.4}{What is data bias in AI?}{1}{L1}{4}
\q{1.5}{What are the two main components of the CNS?}{1}{L1}{3}
\q{1.6}{Which test is commonly used to diagnose epilepsy?}{1}{L1}{3}
\q{1.7}{What type of information does EEG provide that CT and MRI cannot?}{1}{L2}{3}
\q{1.8}{What is the purpose of the 10--20 system in EEG?}{2}{L2}{3}
\end{gtable}
\btfoot{|l|l|l|*{10}{C{0.8cm}|}}{\hline
\multirow{2}{*}{\shortstack{Marks\\Distribution}} & \multicolumn{2}{c|}{Particulars} & CO1 & CO2 & CO3 & CO4 & L1 & L2 & L3 & L4 & L5 & L6\\\cline{2-13}
 & Test & \shortstack{Max\\Marks} & - & - & 05 & 05 & 05 & 05 & - & - & - & -\\\hline}
\par\vspace{6pt}
\begin{gtable}[PART-B]{Sl.\newline No.}{}{M}{BT}{CO}
\q{1.}{Explain the basic functioning of the Central Nervous System (CNS). Describe its main components and their roles.}{10}{L1}{3}
\q{2.}{Explain how CT, MRI, and EEG complement each other in CNS diagnosis.}{10}{L2}{3}
\q{3.a}{List the common ailments of the CNS. Explain any four with their symptoms in brief.}{5}{L1}{3}
\q{3.b}{Discuss the major challenges in deploying AI in healthcare, focusing on data bias, privacy, and integration into clinical workflow.}{5}{L2}{4}
\q{4.}{Explain the concept of Artificial Intelligence (AI), Machine Learning (ML), and Deep Learning (DL) in the context of healthcare}{10}{L1}{4}
\q{5.}{Explain the types of data used in healthcare. Differentiate between structured and unstructured data with examples.}{10}{L2}{4}
\end{gtable}
\btfoot{|l|l|l|*{10}{C{0.8cm}|}}{\hline
\multirow{2}{*}{\shortstack{Marks\\Distribution}} & \multicolumn{2}{c|}{Particulars} & CO1 & CO2 & CO3 & CO4 & L1 & L2 & L3 & L4 & L5 & L6\\\cline{2-13}
 & Test & \shortstack{Max\\Marks} & - & - & 25 & 25 & 25 & 25 & - & - & - & -\\\hline}
\end{iempaperB}

% ---------- CH375TGD Green and Hydrogen Technology (CIE 3, Improvement Test)
\begin{gepaper}
\paperinfo{program={CHEMICAL ENGINEERING},title={GREEN AND HYDROGEN TECHNOLOGY},code={CH375TGD},exam={CONTINUOUS INTERNAL EVALUATION \ \ 3 (Improvement Test)},ggroup={INSTITUTIONAL ELECTIVE: GROUP D},sem={VII},month={Jan 2026},maxmarks={60},gdate={14.01.2026},mtime={13:20 - 15:20},duration={120 Min},
 instr={All the questions are compulsory. No additional booklet will be provided.}}
\gehead
\begin{getable}
\gq{1}{What is hydrodesulfurization?}{02}{01}{03}
\gq{2}{Define hydrogenation?}{02}{01}{02}
\gq{3}{List the applications of hydrogen in fertilizer industry.}{02}{01}{01}
\gq{4}{What is the use of hydrogen in aviation sector?}{02}{01}{04}
\gq{5}{What is the role of hydrogen in fuel cells?}{02}{01}{04}
\gepart{\bfseries PART B}
\gq{1}{How are hydrogen hazards classified? Explain deflagration, detonation, explosion, fireball, and flash fire.}{10}{02}{01}
\gq{2}{List and explain the hydrogen risk components related to material properties and hydrogen handling aspects.}{10}{01}{04}
\gq{3}{List the safety features to be incorporated in electrolyzers, pressure vessels, trailers (gaseous and liquid), pipelines, and fuel cell vehicles for hydrogen handling.}{10}{03}{03}
\gq{4}{Give an explanatory note on the applications of hydrogen in refineries and petrochemicals.}{10}{02}{02}
\gq{5}{Explain the role of hydrogen in transport and automotive sectors.}{10}{02}{02}
\end{getable}
\par\vspace{6pt}\noindent{\fontsize{12}{14}\selectfont Course Outcomes}\par\vspace{2pt}
\noindent\begin{tabular}{|C{0.9cm}|p{15.3cm}|}\hline \multicolumn{2}{|l|}{At the end of the course the student will be able to,}\\\hline
1 & Understand the importance of hydrogen and its use as an energy carrier\\\hline
2 & Explain the production, storage and handling of hydrogen\\\hline
3 & Analyze the need for hydrogen as an alternate fuel and the associated challenges\\\hline
4 & Appraise the importance of safety, regulations and codes\\\hline\end{tabular}
\par\vspace{10pt}
\noindent\begin{tabular}{|C{2.0cm}|*{4}{C{1.2cm}|}*{6}{C{1.1cm}|}}\hline Cos/BTL & CO1 & CO2 & CO3 & CO4 & L1 & L2 & L3 & L4 & L5 & L6\\\hline Marks & 10 & 20 & 10 & 10 & 10 & 30 & 10 & - & - & -\\\hline\end{tabular}
\end{gepaper}

% ---------- AS375TGB Aircraft Systems (Offline Quiz & Test, Improvement)
\begin{xpaper}
\begin{center}{\bfseries\fontsize{13}{15}\selectfont Department of}\\[2pt]{\bfseries\fontsize{16}{19}\selectfont Aerospace Engineering}\end{center}\vspace{-6pt}
\noindent\begingroup\renewcommand{\arraystretch}{1.3}\fontsize{11.5}{13}\selectfont
\begin{tabular}{|p{3.4cm}|p{4.0cm}|p{5.6cm}|p{2.6cm}|}\hline
Date & 14\textsuperscript{th} January \ 2026 & Maximum Marks & 60\\\hline
Course Code & AS375TGB & Duration & 120Min\\\hline
Sem & VII Semester & \multicolumn{2}{l|}{Offline: \ Quiz \&Test ( Improvement)}\\\hline
\multicolumn{4}{|c|}{AIRCRAFT SYSTEMS}\\\hline
\end{tabular}\endgroup\par\vspace{8pt}
\begin{gtable}[PART -A]{Sl No}{Questions}{M}{BT}{CO}
\q{1}{List any two fire suppression agents used in aircraft.}{2}{2}{2}
\q{2}{Why is cabin pressurization important at high altitudes?}{2}{2}{3}
\q{3}{State one advantage of fully powered flight controls over conventional systems.}{2}{3}{1}
\q{4}{What is the purpose of trim tabs in flight control systems?}{2}{4}{1}
\q{5}{Differentiate between primary and secondary flight controls.}{2}{2}{4}
\end{gtable}
\par\vspace{6pt}
\begin{gtable}[PART-B]{Sl. No.}{Questions}{M}{BT}{CO}
\q{1.}{Evaluate the design and operation of aircraft fire detection and suppression systems with respect to safety requirements.}{10}{L4}{4}
\q{2.}{Compare air cycle and vapour cycle air-conditioning systems used in aircraft. Analyze their suitability for different aircraft categories.}{10}{L3}{3}
\q{3.}{Explain the history of airplanes and its influence on the development of modern flight control systems}{10}{L2}{1}
\q{4.}{Analyze how redundancy and fault tolerance are incorporated in electronic flight control systems.}{10}{L3}{2}
\q{5.}{Analyze the limitations of conventional flight control systems and justify the need for power-assisted controls.}{10}{L3}{4}
\end{gtable}
\btfoot{|l|l|l|*{10}{C{0.8cm}|}}{\hline
\multirow{2}{*}{\shortstack{Marks\\Distribution}} & \multicolumn{2}{c|}{Particulars} & CO1 & CO2 & CO3 & CO4 & L1 & L2 & L3 & L4 & L5 & L6\\\cline{2-13}
 & Test & \shortstack{Max\\Marks} & 14 & 12 & 12 & 22 & 00 & 14 & 34 & 12 & 0 & \\\hline}
\end{xpaper}

% ---------- MA375TGR Mathematics of Music (CIE 3)
\begin{xpaper}
\par\vspace{4pt}
\noindent\begingroup\renewcommand{\arraystretch}{1.5}\bfseries\fontsize{11.5}{13}\selectfont
\begin{tabular}{|p{5.0cm}|C{4.6cm}|p{6.5cm}|}\hline
Date: 14-01-2026 & CIE 3 & Max. Marks: 10+50\\\hline
Semester: VII & UG & Duration: 2 Hrs {\footnotesize (1:20PM- 3:20PM)}\\\hline
\multicolumn{2}{|l|}{Course Title: Mathematics of Music} & Course code: MA375TGR\\\hline
\end{tabular}\endgroup\par\vspace{2pt}
\centerline{\bfseries\fontsize{13}{15}\selectfont Department of Mathematics}\par\vspace{4pt}
\begin{gtable}{S No}{Quiz}{M}{BT}{CO}
\q{1.}{Define avartana and sama in the context of the Carnatic tala system.}{2}{1}{2}
\q{2.}{State the tihai formula 3P + 2G = T and explain what each variable represents.}{2}{2}{1}
\q{3.}{What is a fractal? Give one example of a deterministic fractal.}{2}{2}{2}
\q{4.}{Define critical band in psychoacoustics and state its significance for consonance perception.}{2}{3}{3}
\q{5.}{Calculate the total number of aksharas in Misra-jati Rupaka Tala.}{2}{2}{3}
\end{gtable}
\par\vspace{8pt}
\begin{gtable}{S No}{Test}{M}{BT}{CO}
\q{1a}{Explain the hierarchical organization of time in Carnatic music: matra, akshara, and avartana. How do these relate mathematically through jati and nadai?}{5}{2}{1}
\q{1b}{A tala has structure I O I I with Khanda jati. Calculate the number of aksharas in one cycle. If rendered in Misra nadai, how many matras are in one avartana?}{5}{2}{1}
\q{2a}{Design a tihai for 2 cycles of Adi Tala in Chatusra nadai (64 matras). Specify the pattern length P and gap G using the tihai formula 3P + 2G = T. Verify that your design ends at sama. (ii)}{5}{3.}{2}
\q{2b}{A korvai has purvanga of 12 matras and uttaranga of 8 matras. If repeated 3 times, calculate the total length. Does it fit exactly in 2 cycles of Rupaka Tala (Chatusra jati, Chatusra nadai)? Justify.}{5}{3}{2}
\q{3a}{Explain the concept of 1/f noise (pink noise) in music as discovered by Voss and Clarke. How does it relate to entropy and redundancy in musical composition?}{5}{3}{4}
\q{3b}{Calculate the Shannon entropy and redundancy for a four-note melody system with probabilities: Sa (0.4), Ri (0.3), Ga (0.2), Ma (0.1). Interpret your results in terms of listener engagement.}{5}{3}{4}
\q{4}{(i) Define critical bandwidth using Zwicker's formula. Calculate the critical bandwidth at 520 Hz. (ii) Using your result from (i), analyze whether an equal temperament minor third interval (520 Hz and \blank[1cm] Hz) would be perceived as consonant or dissonant. Explain the relationship between critical bands and the perception of roughness.}{\illegc}{\illegc}{\illegc}
\q{5a}{Define a first-order Markov chain. Construct a valid 3$\times$3 transition matrix for generating melodies using notes \{C, D, E\} where stepwise motion has probability 0.5, staying on the same note has probability 0.3, and skipping has probability 0.2. Verify your matrix is row-stochastic.}{\illegc}{\illegc}{\illegc}
\q{5b}{Explain the missing fundamental phenomenon in psychoacoustics. A complex tone contains harmonics at 200 Hz, 300 Hz, 400 Hz, and 500 Hz but lacks the fundamental. What pitch will a listener perceive and why?}{\illegc}{\illegc}{\illegc}
\end{gtable}
\par\vspace{2pt}\centerline{\small\itshape [The M/BT/CO columns of Q4 and Q5 are cut off in the scanned copy.]}
\end{xpaper}

% ---------- ET375TGL Space Technology and Applications (CIE - 3)
\begin{xpaper}
\xdept{Department of Electronics and Telecommunication Engineering}{Date: 14/01/2026}{CIE - 3}{Max. Marks: 60}{Semester: 7}{UG}{Duration: 2 Hrs}{Course Title: Space Technology and Applications}{Course Code:ET375TGL}
\begin{gtable}{Q.No}{QUIZ}{M}{BT}{CO}
\q{1}{List the different types of remote sensing systems based on the range of electromagnetic radiation.}{2}{L2}{CO3}
\q{2}{Compare the active and passive remote sensing systems}{2}{L3}{CO3}
\q{3}{What are the main functions of the space shuttle? List its main components.}{2}{L2}{CO4}
\q{4}{What are the design constraints of the planetary payloads?}{2}{L2}{CO4}
\q{5}{What are the communication challenges faced during lunar missions?}{2}{L2}{CO4}
\qpart{TEST}
\q{1}{Explain the working principle of passive remote sensing systems along with advantages and limitations.}{10}{L2}{CO3}
\q{2}{Explain the following with respect to satellite image\newline i) Spatial Resolution\newline ii) Spectral Resolution\newline iii) Temporal Resolution\newline iv) Radiometric Resolution}{10}{L2}{CO3}
\q{3}{Explain the Image processing techniques to process the satellite data into useful information.}{10}{L2}{CO3}
\q{4}{Illustrate the following applications of remote sensing satellites\newline i) Land Cover Classification\newline ii) Land Cover Change Detection}{10}{L3}{CO3}
\q{5}{Explain the Deep space planetary mission along with challenges and requirements.}{10}{L2}{CO4}
\end{gtable}
\btfoot{|l|l|l|*{10}{C{0.8cm}|}}{\hline
\multirow{2}{*}{\shortstack{Marks\\Distribution}} & \multicolumn{2}{c|}{Particulars} & CO1 & CO2 & CO3 & CO4 & L1 & L2 & L3 & L4 & L5 & L6\\\cline{2-13}
 & Test & Max Marks & - & - & 44 & 16 & - & 48 & 12 & - & - & -\\\hline}
\end{xpaper}

% ---------- IM375TGM Project Management (Institutional Elective), Improvement CIE
% (scan p22 and p23 are two photos of the same paper; p22 is complete)
\begin{iempaperB}
\paperinfo{ayear={2025-2026 (ODD Sem)},date={14.01.2026},code={IM375TGM},sem={VII Semester},exam={Improvement CIE},maxmarks={10 + 50},duration={120 Mins},
 title={PROJECT MANAGEMENT (Institutional Elective)},shownote=0}
\iemheadB
\onenote{Answer all the Questions.}
\begin{gtable}[PART -- A]{Sl.\newline No.}{Questions}{M}{BT}{CO}
\q{1}{Name any two project scheduling techniques.}{02}{L1}{CO1}
\q{2}{Differentiate between CPM and PERT}{02}{L2}{CO1}
\q{3}{List any two positive strategies to mitigate risk.}{02}{L2}{CO2}
\q{4}{Highlight any two quantitative risk analysis tools.}{02}{L1}{CO1}
\q{5}{Mention any two tools used in qualitative risk analysis.}{02}{L2}{CO2}
\end{gtable}
\begin{gtable}[PART -- B]{Sl.\newline No.}{Questions}{M}{BT}{CO}
\q{1.}{Draw the network diagram and find forward pass \& backword pass computations, identify critical path for the following:
\qtabc{|l|c|c|c|c|c|c|c|}{\hline Activity & 1 -- 2 & 1 -- 3 & 1 -- 4 & 2 -- 5 & 3 -- 7 & 4 -- 6 & 5 -- 7\\\hline Duration & 10 & 8 & 9 & 8 & 16 & 7 & 7\\\hline Activity & 5 -- 8 & 6 -- 7 & 6 -- 9 & 7 -- 10 & 8 -- 10 & 9 -- 10 & \multicolumn{1}{c}{}\\\cline{1-7} Duration & 6 & 7 & 5 & 12 & 13 & 15 & \multicolumn{1}{c}{}\\\cline{1-7}}}{15}{L4}{CO4}
\q{2.}{A manufacturing company is launching a new product and needs to plan its project schedule efficiently. As a project manager, explain the steps involved in project implementation scheduling and discuss how effective time management and resource allocation methods can optimize the project timeline.}{10}{L2}{CO2}
\q{3.}{The following network diagram represents activities associated with a project:
\qtabc{|l|*{9}{c|}}{\hline Activity & A & B & C & D & E & F & G & H & I\\\hline Optimistic Time & 5 & 18 & 26 & 16 & 15 & 6 & 7 & 7 & 3\\\hline Pessimistic Time & 10 & 22 & 40 & 20 & 25 & 12 & 12 & 9 & 5\\\hline Most Likely Time & 8 & 20 & 33 & 18 & 20 & 9 & 10 & 8 & 4\\\hline}
\begin{romi}\item[i.] Find expected completion time and variance of each activity\item[ii.] Draw Network diagram\item[iii.] The critical path and project duration\item[iv.] The duration of the project that will have 95 per cent chance of being completed.\end{romi}}{15}{L4}{CO4}
\q{4.}{Explain the application of Project Life Cycle Costing in an enterprise software implementation project.}{10}{L4}{CO4}
\end{gtable}
\stars{*****}
\end{iempaperB}

\end{document}
